\section{Investigación: simulación, algoritmo genético y alternativas}

\subsection{Simulación de eventos discretos}
El estado de un DE cambia en instantes discretos (llegada, inicio y fin de atención), y la
congestión surge de la interacción entre llegadas aleatorias y capacidad limitada. La
simulación de eventos discretos (DES) avanza el reloj solo entre eventos y permite registrar
cada espera y comparar políticas bajo las mismas secuencias aleatorias. Se usó SimPy
\citep{simpy412} para el reloj, la agenda de eventos y los recursos con prioridad; el
equipo desarrolló los procesos de paciente, el cambio de capacidad por bloque, la cola
clínica, las métricas y las verificaciones.

\subsection{Algoritmo genético restringido}
Un algoritmo genético (GA) mantiene una población de soluciones y las mejora mediante
selección, cruce y mutación \citep{eiben2015evolutionary}. Aquí un individuo es un vector
entero $x=(x_0,\dots,x_5)$ con los equipos de cada bloque de cuatro horas. La evaluación
de un individuo es una simulación estocástica, por lo que el GA usa la DES como función de
aptitud ruidosa. Las restricciones se manejan con la regla de factibilidad de
\citet{deb2000constraints}: una solución factible domina a una no factible; entre factibles
gana el menor objetivo; entre no factibles gana la menor violación. La literatura reciente
sobre dotación de personal en emergencias combina simulación y optimización de forma
similar \citep{wang2025staffing}.

\subsection{Alternativas consideradas}
\begin{table}[ht]
\centering\small
\caption{Comparación de técnicas de optimización para el problema de dotación horaria.}
\label{tab:alternativas}
\begin{tabularx}{\linewidth}{@{}l X X@{}}
\toprule
Alternativa & Ventaja & Por qué no es la técnica principal \\
\midrule
Enumeración exhaustiva & Encuentra el mejor valor del espacio evaluado & Cada candidato exige varias simulaciones; el costo crece con límites, bloques y réplicas \\
Programación lineal entera & Ofrece garantías en problemas lineales & La espera y sus percentiles emergen de una simulación no lineal y estocástica \\
Búsqueda aleatoria & Sencilla; útil como referencia & No reutiliza información de soluciones prometedoras \\
Enjambre de partículas & Buen desempeño en espacios continuos & Representación entera, restringida, con reparación frecuente \\
Algoritmo genético & Maneja vectores enteros, operadores propios y restricciones & No garantiza optimalidad global; exige controlar ruido y convergencia \\
\bottomrule
\end{tabularx}
\end{table}

\subsection{Limitaciones de la técnica}
Un GA no demuestra optimalidad global y puede converger prematuramente. Con pocas
réplicas por candidato puede seleccionar por ruido; por eso la candidata se reevalúa con
semillas independientes (sección~\ref{sec:diseno}). La investigación previa del proyecto
planteó una búsqueda aleatoria con el mismo presupuesto de evaluaciones como comparador;
\emph{la evidencia final no incluye ese comparador}, de modo que no se puede afirmar cuánto
aporta la estructura evolutiva frente a probar horarios al azar.
